\documentclass{article}
\usepackage{fullpage}
\usepackage{a4}
\usepackage{amsmath}
\parindent 0pt
\parskip 6pt

%\newcommand{\text}[1]{\mbox{#1}}
\begin{document}
 
\centerline{\bf Conversation between Peter Buneman and  GPT 4o on 27 Oct 2025}




{\bf Me}: I have a square matrix with 1s on the diagonal and non-positive
numbers everywhere else. The sum of each row is strictly less than
1. Prove that this matrix has an inverse

{\bf It}: To prove that such a matrix has an inverse, we'll consider a square matrix \( A = [a_{ij}] \) of order \( n \). The matrix \( A \) is defined such that \( a_{ii} = 1 \) for all \( i \), and \( a_{ij} \leq 0 \) for \( i \neq j \). Additionally, the row sum condition is given by:

\[
\sum_{j=1}^{n} a_{ij} < 1 \quad \text{for all } i
\]

This condition implies:

\[
1 + \sum_{j \neq i} a_{ij} < 1 \quad \Rightarrow \quad \sum_{j \neq i} a_{ij} < 0
\]

To prove that \( A \) is invertible, we'll show that it is diagonally dominant, which is a sufficient condition for invertibility.

A matrix is strictly diagonally dominant if for each row the absolute value of the diagonal entry is greater than the sum of the absolute values of the non-diagonal entries in that row. Formally, \( A \) is strictly diagonally dominant if:

\[
|a_{ii}| > \sum_{j \neq i} |a_{ij}|
\]

Given the conditions, we have on the diagonal \( a_{ii} = 1 \). Off-diagonal elements \( a_{ij} \) are non-positive, hence:

\[
|a_{ij}| = -a_{ij} \quad \text{since} \quad a_{ij} \leq 0
\]

The condition for strict diagonal dominance becomes:

\[
1 > \sum_{j \neq i} (-a_{ij})
\]

However, since \( \sum_{j \neq i} a_{ij} < 0 \), we know:

\[
\sum_{j \neq i} -a_{ij} = -\sum_{j \neq i} a_{ij} > 0
\]

Thus, the diagonal dominance condition \( 1 > -\sum_{j \neq i} a_{ij} \) holds.

The Gershgorin Circle Theorem also supports this conclusion. According to this theorem, every eigenvalue of \( A \) lies within at least one of the Gershgorin discs centered at \( a_{ii} \) with radii \( R_i = \sum_{j \neq i} |a_{ij}| \). Here, all our radii \( R_i \) satisfy \( R_i < 1 \). Hence, each disc is centered at 1 and extends less than 1 away from it, ensuring that 0 is not an eigenvalue.

Therefore, all eigenvalues of \( A \) are non-zero, ensuring the matrix is invertible.

In conclusion, the matrix \( A \) has an inverse because it is
strictly diagonally dominant and satisfies the conditions of the
Gershgorin Circle Theorem which confirm that it is non-singular.


{\bf Me}: Please give me a reference for diagonal dominance being a
sufficient condition for invertibility.

{\bf It}: You can find references to diagonal dominance as a sufficient condition for invertibility in most linear algebra textbooks. One classic reference is:

- **"Matrix Analysis" by Roger A. Horn and Charles R. Johnson.**

This book provides a thorough explanation of diagonal dominance and
its implications for the invertibility of matrices, including
necessary theorems and proofs.

{\bf Me}: Thank you.  Can you also give me a proof of this?

{\bf It}: Certainly! Here's a proof that a strictly diagonally dominant matrix is invertible:

### Definitions
A square matrix \( A = [a_{ij}] \) is strictly diagonally dominant if for every row \( i \):

\[
|a_{ii}| > \sum_{j \neq i} |a_{ij}|
\]

### Proof of Invertibility

Assume \( A \) is strictly diagonally dominant. We want to show that \( A \) is invertible, which means it has no zero eigenvalues.

#### Step 1: Use the Gershgorin Circle Theorem

The Gershgorin Circle Theorem states that every eigenvalue of \( A \) lies within at least one Gershgorin disc. For row \( i \), the disc is centered at \( a_{ii} \) with radius \( R_i = \sum_{j \neq i} |a_{ij}| \).

Strict diagonal dominance implies:

\[
|a_{ii}| > R_i
\]

So, the disc is entirely within the range that prevents zero from being an eigenvalue. Specifically, since:

\[
|a_{ii}| - R_i > 0
\]

every eigenvalue \( \lambda \) must satisfy:

\[
|\lambda - a_{ii}| < R_i
\]

This condition implies that \( \lambda \neq 0 \) because \( 0 \) would not satisfy the inequality given the strict dominance by \( |a_{ii}| \).

#### Step 2: Conclusion

Since none of the Gershgorin discs include zero, \( A \) cannot have zero as an eigenvalue. 

#### Step 3: Result

Because all eigenvalues of \( A \) are non-zero, \( A \) is invertible.

This concludes the proof that strict diagonal dominance implies
invertibility.

{\bf Me}: Thank you. Can you describe the Gershgorin disc  In particular you
say it is centered at \(a_{ii}\), but this is a value, it is not the
coordinates of a point.

{\bf It}: Certainly! In the context of complex numbers and eigenvalues, the Gershgorin disc for each row \( i \) of a matrix \( A \) is defined in the complex plane.

### Description of a Gershgorin Disc

- **Center**: The center of the disc is the complex number \( a_{ii} \), which represents a specific point on the complex plane (a real number in this context, if the matrix has real values).

- **Radius**: The radius of the disc is given by the sum of the absolute values of the non-diagonal entries in the same row:
  \[
  R_i = \sum_{j \neq i} |a_{ij}|
  \]

### Interpretation

For each eigenvalue \( \lambda \) of the matrix, the Gershgorin Circle Theorem states:
\[
|\lambda - a_{ii}| < R_i
\]

This means that the eigenvalue \( \lambda \) lies within or on the
boundary of this disc in the complex plane. Each Gershgorin disc
encompasses potential eigenvalues, providing useful bounds and
insights about their locations.

{\bf Me}: I see.   Can you also give me a proof of the Gershgorin Circle
Theorem

{\bf It}:
Certainly! Here's a proof of the Gershgorin Circle Theorem:

### Gershgorin Circle Theorem

For an \( n \times n \) matrix \( A = [a_{ij}] \), every eigenvalue \( \lambda \) of \( A \) satisfies:

\[
|\lambda - a_{ii}| \leq \sum_{j \neq i} |a_{ij}|
\]

for at least one \( i \).

### Proof

1. **Eigenvalue Definition**:
   Suppose \( \lambda \) is an eigenvalue of \( A \) with eigenvector \( \mathbf{v} = [v_1, v_2, \ldots, v_n]^T \), where \( \mathbf{v} \neq \mathbf{0} \).

2. **Matrix Equation**:
   \( A\mathbf{v} = \lambda\mathbf{v} \).

3. **Rewrite the Equation**:
   For entry \( i \), this becomes:
   \[
   \sum_{j=1}^n a_{ij} v_j = \lambda v_i
   \]

4. **Rearrange**:
   \[
   (\lambda - a_{ii})v_i = \sum_{j \neq i} a_{ij} v_j
   \]

5. **Take Absolute Values**:
   \[
   |\lambda - a_{ii}||v_i| = \left|\sum_{j \neq i} a_{ij} v_j\right|
   \]

6. **Use Triangle Inequality**:
   \[
   \left|\sum_{j \neq i} a_{ij} v_j\right| \leq \sum_{j \neq i} |a_{ij}||v_j|
   \]

7. **Divide by \(|v_i|\) (assuming \(v_i \neq 0\))**:
   \[
   |\lambda - a_{ii}| \leq \sum_{j \neq i} |a_{ij}|\left|\frac{v_j}{v_i}\right|
   \]

8. **Select Maximum Value**:
   Assume \(|v_k| = \max \{|v_1|, |v_2|, \ldots, |v_n|\}\).

9. **Estimate**:
   If \(i = k\), then:
   \[
   |\lambda - a_{ii}| \leq \sum_{j \neq i} |a_{ij}|\frac{|v_j|}{|v_i|} \leq \sum_{j \neq i} |a_{ij}|
   \]

10. **Conclusion**:
    There must exist such an \( i \) where the inequality holds, ensuring every eigenvalue is within at least one Gershgorin disc.

This proves the theorem, showing every eigenvalue lies within one of
the discs centered at diagonal entries \( a_{ii} \) with radii defined
by the sums of the non-diagonal entries.

{\bf Me}: Thank you. That was helpful

It: You're welcome! If you have any more questions, feel free to ask.

\end{document}
